\documentclass[11pt]{scrartcl}
\usepackage[a4paper,margin=2.6cm]{geometry}
\usepackage{../common/rlstudy}

\title{Extended reading --- Week 19: DPO and the Preference Zoo}
\date{\today}
\begin{document}
\maketitle

\section*{How to read this list}
Each paper below: full citation, one-paragraph summary, a
\emph{Why it matters for this week} note, and one annotated key
passage with brief commentary.

\section*{1. Rafailov et al.\ 2023 --- Direct Preference Optimization}
\textbf{Citation.} R. Rafailov, A. Sharma, E. Mitchell, S. Ermon,
C.\,D.\ Manning, C. Finn, \emph{Direct Preference Optimization: Your
Language Model is Secretly a Reward Model}, arXiv:2305.18290, NeurIPS
2023.

\paragraph{Summary.} Derives, in three pages, the equivalence between
the KL-regularised RL objective and a closed-form supervised
contrastive loss. The reward model can be re-parametrised through
the policy as $r_\theta = \beta\log(\pi_\theta/\pi_\text{ref})$; the
partition function cancels in pairwise comparisons; what remains is
a logistic-regression-style loss on the policy. Empirically, DPO
matches or beats PPO-RLHF on summarisation and dialogue at a
fraction of the compute and engineering cost.

\paragraph{Why it matters for this week.} This is the paper the
whole week is built around. The 4-step derivation in our notes is
their Theorem~1.

\paragraph{Annotated key passage.}
\begin{quote}\itshape
``This change of variables enables us to express the RL objective in
terms of the policy directly, avoiding the need to learn an
explicit reward model... The resulting optimization problem can be
solved using a standard supervised learning algorithm, without the
need for RL.''
\end{quote}
\noindent The two clauses are operationally massive. ``Avoiding the
need to learn an explicit reward model'' kills three of the four
models in the RLHF pipeline. ``Without the need for RL'' kills the
remaining engineering surface (rollouts, GAE, clipping, KL
controller, value head). The community's adoption of DPO within
twelve months was almost entirely a story about this engineering
simplification.

\section*{2. Azar et al.\ 2024 --- IPO and the $\Psi$PO framework}
\textbf{Citation.} M.\,G.\ Azar, M. Rowland, B. Piot, D. Guo, D.
Calandriello, M. Valko, R. Munos, \emph{A General Theoretical
Paradigm to Understand Learning from Human Preferences},
arXiv:2310.12036, AISTATS 2024.

\paragraph{Summary.} Generalises DPO into $\Psi$PO, a family
parametrised by the choice of link function $\Psi$. Shows DPO's
logistic link saturates as preferences become deterministic,
collapsing the implicit KL constraint. Proposes IPO --- a squared
loss, $\Psi(u) = (u - 1/(2\beta))^2$ --- whose unique optimum has a
\emph{finite} margin and avoids the saturation pathology.

\paragraph{Why it matters for this week.} The link-function plot in
the notes (\texttt{fig\_links.pdf}) is theirs.
IPO is the cleanest theoretical advance over DPO; KTO, ORPO, and
SimPO are best understood as alternative \emph{structural} (not
link-function) choices on top of this framework.

\paragraph{Annotated key passage.}
\begin{quote}\itshape
``As the preferences become more deterministic, the empirical
optima of DPO drive the implicit reward gap to infinity, regardless
of the value of $\beta$. IPO instead has a finite optimum...''
\end{quote}
\noindent The two-sentence rebuke of DPO. The IPO ``finite optimum''
is also the reason IPO does \emph{not} need a KL leash: the squared
loss provides its own regularisation by pinning the margin to a
finite target.

\section*{3. Ethayarajh et al.\ 2024 --- KTO}
\textbf{Citation.} K. Ethayarajh, W. Xu, N. Muennighoff, D. Jurafsky,
D. Kiela, \emph{KTO: Model Alignment as Prospect Theoretic
Optimization}, arXiv:2402.01306, 2024.

\paragraph{Summary.} Frames alignment through Kahneman--Tversky
prospect theory: human preferences are noisy and asymmetric --- gains
relative to a reference are felt less strongly than equal-sized
losses. KTO replaces the BT pairwise model with a per-sample
\emph{utility} that is concave on gains and convex on losses,
allowing training on \emph{unpaired} thumbs-up / thumbs-down labels.

\paragraph{Why it matters for this week.} Production data is rarely
in clean pairwise form; KTO is the practical way to use the messy,
unpaired feedback that the deployed product accumulates. The
synthetic insurance matching story in the notes lives here.

\paragraph{Annotated key passage.}
\begin{quote}\itshape
``KTO does not need preferences --- only a binary signal of whether
an output is desirable or undesirable for a given input. This makes
it far easier to use in the real world, where this kind of data is
much more abundant.''
\end{quote}
\noindent The pragmatic case for KTO in one sentence. The deeper
theoretical contribution is that the prospect-theory utility is
\emph{not} an arbitrary heuristic: it falls out of empirically
calibrated human-decision data, giving KTO a behavioural-science
foundation that DPO lacks.

\section*{4. Hong et al.\ 2024 --- ORPO}
\textbf{Citation.} J. Hong, N. Lee, J. Thorne, \emph{ORPO: Monolithic
Preference Optimization without Reference Model}, arXiv:2403.07691,
EMNLP 2024.

\paragraph{Summary.} A single-stage method that fuses SFT NLL with a
log-odds-ratio penalty on preference pairs, eliminating the
reference policy entirely. The result: a one-stage pipeline that
trains on the same data DPO uses, at half the memory, with a small
quality bump on AlpacaEval and IFEval.

\paragraph{Why it matters for this week.} ORPO is the natural
contrast to DPO on the ``do I need a reference?'' axis. It also
forces us to grapple with numerically stable log-odds-ratio
computation --- a subtle bit of practical engineering exposed in
\texttt{losses.\_log1mexp}.

\paragraph{Annotated key passage.}
\begin{quote}\itshape
``Despite ORPO being a single-step training method, models trained
with ORPO surpass much larger Llama-2-Chat and Zephyr variants
trained with the conventional multi-stage RLHF pipeline.''
\end{quote}
\noindent The bold empirical claim. With caveats: ORPO's
single-stage nature couples preference data quality to SFT data
quality very tightly. On clean curated data it shines; on messy
production data it can amplify bias.

\section*{5. Meng, Xia, Chen 2024 --- SimPO}
\textbf{Citation.} Y. Meng, M. Xia, D. Chen, \emph{SimPO: Simple
Preference Optimization with a Reference-Free Reward},
arXiv:2405.14734, NeurIPS 2024.

\paragraph{Summary.} The most ergonomic DPO variant: reference-free,
length-normalised, and with a target margin parameter $\gamma$. The
implicit reward is the average per-token log-probability. Gemma-2-9B
+ SimPO hits 72.4\% length-controlled win rate on AlpacaEval-2
without RLHF.

\paragraph{Why it matters for this week.} SimPO is the recommended
starting baseline for DPO-family work in 2026: it requires no
reference model, has built-in length-normalisation (which by itself
solves one of DPO's biggest headaches), and has only two
hyper-parameters.

\paragraph{Annotated key passage.}
\begin{quote}\itshape
``We propose SimPO, a simple yet effective approach to preference
optimization that uses the average log probability of a sequence as
the implicit reward. This eliminates the need for a reference model
and naturally aligns with the metric used at inference time.''
\end{quote}
\noindent The ``aligns with inference'' point is subtle and
important. DPO's implicit reward is the \emph{sum} of token
log-probabilities, but decoding ranks candidates by the
length-normalised score (perplexity). SimPO simply aligns training
with the inference objective.

\section*{6. Liu et al.\ 2024 --- Iterative Length-Regularized DPO}
\textbf{Citation.} W. Liu, K. Bai, Y. Ren, Y. Bian, J. Bian, Z.
Wang, \emph{Iterative Length-Regularized Direct Preference
Optimization: A Case Study on Improving 7B Language Models to
GPT-4 Level}, arXiv:2406.11817, 2024.

\paragraph{Summary.} Two ideas: (i) iterate DPO by regenerating
preferences with the latest policy, refreshing the implicit-reward
target; (ii) add a length-regularization term to the DPO loss to
combat length inflation. The combination pushes Llama-3-8B to
GPT-4-level on AlpacaEval, a striking result given the model size.

\paragraph{Why it matters for this week.} Bridges DPO and the
iterative-DPO / self-rewarding literature of week~20. Also the most
direct piece of evidence that length bias --- the failure mode SimPO
addresses architecturally --- can be fixed equally well at the loss
level.

\paragraph{Annotated key passage.}
\begin{quote}\itshape
``Iterative DPO without length regularization can win the
AlpacaEval-2 game on raw win rate while losing the length-controlled
win rate by a wide margin. This is a clear signal that the policy
is being rewarded for verbosity rather than quality.''
\end{quote}
\noindent The honest empirical observation that AlpacaEval-2's
length-controlled metric was \emph{designed} to expose: optimisers
hunt the easiest path to winning, and verbosity is the easiest path
under sum-log-prob accounting.

\newpage
\section*{7. Xu et al.\ 2024 --- CPO for machine translation}
\textbf{Citation.} H. Xu, A. Sharaf, Y. Chen, W. Tan, L. Shen,
B. Van Durme, K. Murray, Y.\,J.\ Kim, \emph{Contrastive Preference
Optimization: Pushing the Boundaries of LLM Performance in Machine
Translation}, arXiv:2401.08417, 2024.

\paragraph{Summary.} Introduces CPO --- SFT NLL on the chosen
translation plus a contrastive log-ratio between chosen and rejected.
No reference model, halving the memory footprint vs.\ DPO. Achieves
state-of-the-art on multiple WMT language pairs at the time of
publication.

\paragraph{Why it matters for this week.} The clearest empirical
case for the \emph{halfway-house} family (SFT + contrastive, no
reference): MT is a domain where the SFT signal dominates and the
preference signal is polish.

\paragraph{Annotated key passage.}
\begin{quote}\itshape
``By removing the reference model, CPO not only reduces memory
requirements but also liberates the optimisation from the
constraint of having to stay close to a possibly suboptimal
reference.''
\end{quote}
\noindent The two-edged blade of reference-free training: cheaper,
but you lose the regularising effect of the SFT anchor. CPO's
solution --- include the SFT NLL term --- is one of three reasonable
substitutes (the others being a target margin like SimPO's $\gamma$
and a length normaliser).

\section*{8. Yuan et al.\ 2024 --- Self-Rewarding Language Models}
\textbf{Citation.} W. Yuan, R.\,Y.\ Pang, K. Cho, X. Li, S. Sukhbaatar,
J. Xu, J. Weston, \emph{Self-Rewarding Language Models},
arXiv:2401.10020, 2024.

\paragraph{Summary.} The same LLM plays generator, judge, and
learner. LLM-as-a-judge prompting produces preference scores; DPO
updates the policy on those scores; the loop iterates. Three
iterations of self-rewarding Llama-2-70B beats Claude 2, Gemini Pro
and GPT-4-0613 on AlpacaEval-2.

\paragraph{Why it matters for this week.} Bridges DPO and the
iterative-DPO family of week~20. The DPO loss is unchanged; what
changes is the preference data source. The result is striking but
the failure mode --- self-confirmation drift --- is also striking.

\paragraph{Annotated key passage.}
\begin{quote}\itshape
``To use the model itself as a reward judge, we use prompts that
instruct it to assess responses, providing a numerical score that
we can use as a reward signal for further training. Crucially, the
model's ability to judge improves alongside its ability to
generate.''
\end{quote}
\noindent The optimistic claim. The pessimistic one is that the
model's biases as a judge co-evolve with its biases as a generator,
producing a runaway feedback loop. Both claims are true to some
extent; the equilibrium depends on data freshness and the
quality of the initial judge.

\newpage
\section*{Pointers for further depth}

\begin{itemize}[leftmargin=*]
  \item \textsc{Rainbow PO} --- arXiv:2410.04203, 2024. A unifying
        survey that places DPO, IPO, KTO, ORPO, SimPO, CPO inside a
        single ``link function $\times$ normalisation'' lattice.
  \item \textsc{DPO survey} --- arXiv:2410.15595, 2024--2025. A
        comprehensive survey of the DPO family up to late 2024.
  \item \textsc{Tang et al.\ 2024} --- \emph{Generalised Preference
        Optimization}, arXiv:2402.05749. An importance-weighted
        analysis of when DPO and explicit-RM-RL gradients agree.
  \item \textsc{Park et al.\ 2024} --- \emph{Disentangling Length
        from Quality in DPO}, arXiv:2403.19159. A clean empirical
        decomposition of why DPO inflates length.
  \item \textsc{Iterative DPO Study} --- arXiv:2503.12854, 2025.
        Comprehensive study of iterative-DPO design choices: how
        often to refresh, how many pairs per iteration, when to
        stop.
  \item \textsc{Tülu-3} --- arXiv:2411.15124, 2024. The clearest
        recipe for combining SFT, DPO, and RLVR end-to-end.
\end{itemize}

\section*{Sourced from}
\texttt{research/llm\_rl\_sota.md} §2 (the DPO zoo) and §6 (iterative /
online DPO).

\end{document}
